% Un transparent beamer rangé hors d'un dossier slides/ (cas réel du corpus
% S4.0 : td/td4/ du cours d'automatique). P3 ne s'applique pas aux
% transparents (slides.md#sl4) : depuis S5.4, l'outil les reconnaît à leur
% classe (supports du vocabulaire du template), plus seulement à leur dossier.
\documentclass[9pt,t]{beamer}
\begin{document}
\begin{frame}{Exemple}
Pour étendre l'intégrale aux fonctions positives, on pose
\begin{definition}
  Énoncé.
\end{definition}
\end{frame}
\end{document}
